\documentclass[10pt,a4paper]{amsart}
\usepackage{amsmath,amssymb,amsthm,graphicx}
\usepackage{verbatim}
\newenvironment{reference}{\comment}{\endcomment}
%\newenvironment{reference}{}{}
\newenvironment{slogan}{\comment}{\endcomment}
\newenvironment{history}{\comment}{\endcomment}

% For commutative diagrams we use Xy-pic
\usepackage[all]{xy}

% We use 2cell for 2-commutative diagrams.
\xyoption{2cell}
% Theorem environments.
%
\theoremstyle{plain}
\newtheorem{theorem}[subsection]{Theorem}
\newtheorem{proposition}[subsection]{Proposition}
\newtheorem{lemma}[subsection]{Lemma}

\theoremstyle{definition}
\newtheorem{definition}[subsection]{Definition}
\newtheorem{example}[subsection]{Example}
\newtheorem{exercise}[subsection]{Exercise}
\newtheorem{situation}[subsection]{Situation}

\theoremstyle{remark}
\newtheorem{remark}[subsection]{Remark}
\newtheorem{remarks}[subsection]{Remarks}

\numberwithin{equation}{subsection}

% Macros
%
\def\lim{\mathop{\mathrm{lim}}\nolimits}
\def\colim{\mathop{\mathrm{colim}}\nolimits}
\def\Spec{\mathop{\mathrm{Spec}}}
\def\Hom{\mathop{\mathrm{Hom}}\nolimits}
\def\Ext{\mathop{\mathrm{Ext}}\nolimits}
\def\SheafHom{\mathop{\mathcal{H}\!\mathit{om}}\nolimits}
\def\SheafExt{\mathop{\mathcal{E}\!\mathit{xt}}\nolimits}
\def\Sch{\mathit{Sch}}
\def\Mor{\mathop{\mathrm{Mor}}\nolimits}
\def\Ob{\mathop{\mathrm{Ob}}\nolimits}
\def\Sh{\mathop{\mathit{Sh}}\nolimits}
\def\NL{\mathop{N\!L}\nolimits}
\def\CH{\mathop{\mathrm{CH}}\nolimits}
\def\proetale{{pro\text{-}\acute{e}tale}}
\def\etale{{\acute{e}tale}}
\def\QCoh{\mathit{QCoh}}
\def\Ker{\mathop{\mathrm{Ker}}}
\def\Im{\mathop{\mathrm{Im}}}
\def\Coker{\mathop{\mathrm{Coker}}}
\def\Coim{\mathop{\mathrm{Coim}}}

% Boxtimes
%
\DeclareMathSymbol{\boxtimes}{\mathbin}{AMSa}{"02}

%
% Macros for moduli stacks/spaces
%
\def\QCohstack{\mathcal{QC}\!\mathit{oh}}
\def\Cohstack{\mathcal{C}\!\mathit{oh}}
\def\Spacesstack{\mathcal{S}\!\mathit{paces}}
\def\Quotfunctor{\mathrm{Quot}}
\def\Hilbfunctor{\mathrm{Hilb}}
\def\Curvesstack{\mathcal{C}\!\mathit{urves}}
\def\Polarizedstack{\mathcal{P}\!\mathit{olarized}}
\def\Complexesstack{\mathcal{C}\!\mathit{omplexes}}
% \Pic is the operator that assigns to X its picard group, usage \Pic(X)
% \Picardstack_{X/B} denotes the Picard stack of X over B
% \Picardfunctor_{X/B} denotes the Picard functor of X over B
\def\Pic{\mathop{\mathrm{Pic}}\nolimits}
\def\Picardstack{\mathcal{P}\!\mathit{ic}}
\def\Picardfunctor{\mathrm{Pic}}
\def\Deformationcategory{\mathcal{D}\!\mathit{ef}}
\usepackage{cleveref}
\crefname{originalequation}{original equation}{original equations}
\Crefname{originalequation}{Original equation}{Original equations}
\makeatletter
\newlabel{equation-base-change-exact-support}{{0E2K}{source}{}{sourcepointer-1}{}}
\newlabel{equation-base-change-exact-support@cref}{{[originalequation][0][]0E2K}{[1][0][]source}{}{}{}}
\makeatother
\crefname{originallemma}{original lemma}{original lemmas}
\Crefname{originallemma}{Original lemma}{Original lemmas}
\makeatletter
\newlabel{lemma-check-base-change-is-iso}{{0E2L}{source}{}{sourcepointer-2}{}}
\newlabel{lemma-check-base-change-is-iso@cref}{{[originallemma][0][]0E2L}{[1][0][]source}{}{}{}}
\makeatother
\crefname{originallemma}{original lemma}{original lemmas}
\Crefname{originallemma}{Original lemma}{Original lemmas}
\makeatletter
\newlabel{perfect-lemma-quasi-coherence-internal-hom}{{0A6H}{source}{}{sourcepointer-3}{}}
\newlabel{perfect-lemma-quasi-coherence-internal-hom@cref}{{[originallemma][0][]0A6H}{[1][0][]source}{}{}{}}
\makeatother
\crefname{originallemma}{original lemma}{original lemmas}
\Crefname{originallemma}{Original lemma}{Original lemmas}
\makeatletter
\newlabel{more-algebra-lemma-internal-hom-evaluate-tensor-isomorphism}{{0ATK}{source}{}{sourcepointer-4}{}}
\newlabel{more-algebra-lemma-internal-hom-evaluate-tensor-isomorphism@cref}{{[originallemma][0][]0ATK}{[1][0][]source}{}{}{}}
\makeatother
\crefname{originalsection}{original section}{original sections}
\Crefname{originalsection}{Original section}{Original sections}
\makeatletter
\newlabel{dualizing-section-base-change-trivial-duality}{{0E28}{source}{}{sourcepointer-5}{}}
\newlabel{dualizing-section-base-change-trivial-duality@cref}{{[originalsection][0][]0E28}{[1][0][]source}{}{}{}}
\makeatother
\crefname{originalequation}{original equation}{original equations}
\Crefname{originalequation}{Original equation}{Original equations}
\makeatletter
\newlabel{dualizing-equation-base-change}{{0E29}{source}{}{sourcepointer-6}{}}
\newlabel{dualizing-equation-base-change@cref}{{[originalequation][0][]0E29}{[1][0][]source}{}{}{}}
\makeatother
\crefname{originallemma}{original lemma}{original lemmas}
\Crefname{originallemma}{Original lemma}{Original lemmas}
\makeatletter
\newlabel{dualizing-lemma-flat-bc-surjection}{{0BZM}{source}{}{sourcepointer-7}{}}
\newlabel{dualizing-lemma-flat-bc-surjection@cref}{{[originallemma][0][]0BZM}{[1][0][]source}{}{}{}}
\makeatother
\title{The base change map for the sheaf with exact support is an isomorphism for a flat morphism and a pseudo-coherent closed immersion}
\author{The Stacks Project Authors}
\date{stacks-project:duality.tex:local-snapshot-20261012; GNU FDL 1.2 or later}
\begin{document}
\maketitle
\noindent\textit{Editorial note.} All delivered mathematical text is the author's own, copied from the pinned Stacks Project chapter file duality.tex; no statement or proof step was rewritten, reordered or removed.\par\medskip
\noindent\textit{Editorial note.} Preamble provenance: the pinned chapter file opens with an input of the shared preamble and carries no preamble text of its own, so the project's single Stacks-wide preamble staged at sources/stacks-duality/preamble.tex is excerpted verbatim under the recipe fixed by LANE\_BRIEF 22.1, as three spans, lines 8-13 (verbatim package and the reference/slogan/history environments), lines 14-18 (Xy-pic setup) and lines 154-217 (the native theorem declarations and mathematical macros). Nothing was added, removed or redefined; the wrapper class and the amsmath/amssymb/amsthm/graphicx base packages are supplied by the assembler.\par\medskip
\noindent\textit{Editorial note.} The counted statement and proof contain no cross reference, no citation and no included file, so the delivered body is self contained.\par\medskip
\noindent\textit{Editorial note.} Canonical location of the counted proposition: Stacks tag 0E2M (sources/stacks-duality/pinned-original-source.tex line 2595).\par\medskip
\noindent\textit{Allowed standard prerequisites (editorial).}\begin{itemize}
\item The base change map for the sheaf with exact support and the criterion for its being an isomorphism. Same chapter, equation labelled equation-base-change-exact-support, invoked in the statement, and Lemma labelled lemma-check-base-change-is-iso, cited at proof lines 2605-2608. Fixes the arrow and reduces the affine case to an algebraic equality.
\item Comparison of derived Hom between rings and schemes for a pseudo-coherent module. Derived Categories of Schemes, Lemma labelled perfect-lemma-quasi-coherence-internal-hom, cited at proof lines 2608-2614. Justifies checking the equality of ring level Homs.
\item The tensor evaluation isomorphism in the flat case. More on Algebra, Lemma labelled more-algebra-lemma-internal-hom-evaluate-tensor-isomorphism part (3), cited at proof lines 2616-2618, with Dualizing Complexes, Section labelled dualizing-section-base-change-trivial-duality. Supplies the displayed equality.
\item The stalk level description of the base change map and its surjectivity criterion. Dualizing Complexes, Equation labelled dualizing-equation-base-change and Lemma labelled dualizing-lemma-flat-bc-surjection, cited at proof lines 2624-2638. Give the second proof.
\end{itemize}\medskip
\ifdefined\hypertarget\hypertarget{sourcepointer-1}{}\fi
\noindent\textit{Source reference.} Original equation 0E2K: sources/stacks-duality/pinned-original-source.tex, source lines 2533--2533. Same-chapter pointer into the author environment equation at source line 2533; invoked at proof line 2627 of the counted proof. It is retained only as a navigation pointer; every load-bearing step is inside the delivered span.\par\medskip
\ifdefined\hypertarget\hypertarget{sourcepointer-2}{}\fi
\noindent\textit{Source reference.} Original lemma 0E2L: sources/stacks-duality/pinned-original-source.tex, source lines 2561--2561. Same-chapter pointer into the author environment lemma at source line 2561; invoked at proof line 2605 of the counted proof. It is retained only as a navigation pointer; every load-bearing step is inside the delivered span.\par\medskip
\ifdefined\hypertarget\hypertarget{sourcepointer-3}{}\fi
\noindent\textit{Source reference.} Original lemma 0A6H: sources/stacks-chow/tags.txt, source lines 12845--12845. Cross-chapter standard prerequisite located in the project tag list staged in the local snapshot; invoked at proof line 2613 of the counted proof. It is retained only as a navigation pointer; every load-bearing step is inside the delivered span.\par\medskip
\ifdefined\hypertarget\hypertarget{sourcepointer-4}{}\fi
\noindent\textit{Source reference.} Original lemma 0ATK: sources/stacks-chow/tags.txt, source lines 13618--13618. Cross-chapter standard prerequisite located in the project tag list staged in the local snapshot; invoked at proof line 2619 of the counted proof. It is retained only as a navigation pointer; every load-bearing step is inside the delivered span.\par\medskip
\ifdefined\hypertarget\hypertarget{sourcepointer-5}{}\fi
\noindent\textit{Source reference.} Original section 0E28: sources/stacks-chow/tags.txt, source lines 17598--17598. Cross-chapter standard prerequisite located in the project tag list staged in the local snapshot; invoked at proof line 2621 of the counted proof. It is retained only as a navigation pointer; every load-bearing step is inside the delivered span.\par\medskip
\ifdefined\hypertarget\hypertarget{sourcepointer-6}{}\fi
\noindent\textit{Source reference.} Original equation 0E29: sources/stacks-chow/tags.txt, source lines 17599--17599. Cross-chapter standard prerequisite located in the project tag list staged in the local snapshot; invoked at proof line 2636 of the counted proof. It is retained only as a navigation pointer; every load-bearing step is inside the delivered span.\par\medskip
\ifdefined\hypertarget\hypertarget{sourcepointer-7}{}\fi
\noindent\textit{Source reference.} Original lemma 0BZM: sources/stacks-chow/tags.txt, source lines 15055--15055. Cross-chapter standard prerequisite located in the project tag list staged in the local snapshot; invoked at proof line 2638 of the counted proof. It is retained only as a navigation pointer; every load-bearing step is inside the delivered span.\par\medskip
\begin{lemma}
\label{lemma-flat-bc-sheaf-with-exact-support}
In the situation above, assume $f$ is flat and $i$ pseudo-coherent.
Then (\ref{equation-base-change-exact-support}) is an isomorphism
for $K$ in $D^+_\QCoh(\mathcal{O}_X)$.
\end{lemma}

\begin{proof}
First proof. To prove this map is an isomorphism, we may work locally.
Hence we may assume $X$, $X'$, $Z$, $Z'$ are affine, say corresponding
to the rings $A$, $A'$, $B$, $B'$. Then $B$ and $A'$ are tor independent
over $A$. By Lemma \ref{lemma-check-base-change-is-iso} it suffices
to check that
$$
R\Hom_A(B, K) \otimes_A^\mathbf{L} A' =
R\Hom_{A'}(B', K \otimes_A^\mathbf{L} A')
$$
in $D(A')$ for all $K \in D^+(A)$. Here we use
Derived Categories of Schemes, Lemma
\ref{perfect-lemma-quasi-coherence-internal-hom}
and the fact that $B$, resp.\ $B'$ is pseudo-coherent as an
$A$-module, resp.\ $A'$-module
to compare derived hom on the level of rings and schemes.
The displayed equality follows from
More on Algebra, Lemma
\ref{more-algebra-lemma-internal-hom-evaluate-tensor-isomorphism}
part (3). See also the discussion in Dualizing Complexes, Section
\ref{dualizing-section-base-change-trivial-duality}.

\medskip\noindent
Second proof\footnote{This proof shows
it suffices to assume $K$ is in $D^+(\mathcal{O}_X)$.}.
Let $z' \in Z'$ with image $z \in Z$.
First show that (\ref{equation-base-change-exact-support})
on stalks at $z'$ induces the map
$$
R\Hom(\mathcal{O}_{Z, z}, K_z)
\otimes_{\mathcal{O}_{Z, x}}^\mathbf{L} \mathcal{O}_{Z', z'}
\longrightarrow
R\Hom(\mathcal{O}_{Z', z'},
K_z \otimes_{\mathcal{O}_{X, z}}^\mathbf{L} \mathcal{O}_{X', z'})
$$
from Dualizing Complexes, Equation (\ref{dualizing-equation-base-change}).
Namely, the constructions of these maps are identical.
Then apply Dualizing Complexes, Lemma \ref{dualizing-lemma-flat-bc-surjection}.
\end{proof}

\end{document}
